\chapter{Conclusion and Future Work}

\section{Summary of Research}

This thesis investigated the problem of reliable entanglement distribution
in dynamic quantum networks where quantum links experience temporal
degradation, fidelity reduction, and limited resource availability.

Unlike traditional static routing approaches that treat quantum links as
unchanging resources, this work proposed a Temporal-State-Aware Routing
framework that considers the evolution of entanglement states over time.
The proposed framework integrates temporal state evaluation, adaptive route
selection, coherence-budget management, and regeneration mechanisms to
maintain reliable communication paths.

A simulation framework was developed to evaluate the proposed method under
different network conditions. The evaluation considered baseline comparison,
multi-seed experiments, scalability analysis, coherence-time variation,
traffic-load variation, and additional sensitivity studies.

The experimental results demonstrated that considering temporal link
conditions provides more reliable routing decisions compared with static
routing approaches. The proposed framework improves the capability of quantum
networks to maintain usable entanglement paths under dynamic operating
conditions.


\section{Main Contributions}

The main contributions of this thesis are summarized as follows:

\begin{itemize}

\item
A temporal-state-aware quantum routing framework was proposed for
fidelity-preserving entanglement distribution in dynamic quantum networks.
The framework considers the time-dependent evolution of entanglement
resources instead of relying only on static topology information.

\item
A routing strategy was developed that incorporates temporal link conditions,
including link aging, fidelity degradation, and remaining coherence budget,
into route selection decisions.

\item
A regeneration-assisted resource management mechanism was introduced to
recover degraded entanglement resources and improve long-term route
availability.

\item
A simulator was implemented to evaluate routing under different network sizes,
coherence times, traffic conditions, and random seeds.

\item
A comprehensive evaluation framework was established using metrics including
request success ratio, fidelity performance, temporal disconnection ratio,
completion delay, regeneration overhead, and quantum-memory utilization.

\end{itemize}


\section{Limitations}

Although the proposed framework improves routing reliability in simulation,
several limitations remain.

First, the current evaluation is based on simulated quantum network
environments. Real quantum communication systems may introduce additional
hardware-dependent imperfections, noise sources, and operational constraints.

Second, the current model uses a simplified temporal fidelity degradation
process. More detailed physical noise models, including hardware-specific
errors and complex quantum channel behaviors, can be incorporated in future
work.

Third, the regeneration mechanism assumes that recovery resources are
available when required. Practical quantum networks may face limitations in
entanglement generation capability, memory availability, and operational
cost.

These limitations indicate several possible directions for extending the
proposed framework toward more realistic quantum network implementations.


\section{Future Research Directions}

Several research directions can be explored based on this work.

First, future studies can incorporate more realistic quantum noise models and
hardware parameters to improve the physical accuracy of the simulation
framework.

Second, machine learning and reinforcement learning approaches can be
investigated for adaptive routing decisions in larger and more complex
quantum networks.

Third, future work can consider distributed routing architectures where
multiple quantum network nodes collaboratively exchange temporal state
information.

Finally, the proposed framework can be validated using quantum network
simulation platforms and experimental quantum communication testbeds to
evaluate its practical applicability.


\section{Conclusion}

This thesis developed a temporal-state-aware framework for reliable
entanglement distribution in dynamic quantum networks. It combines temporal
resource tracking, adaptive routing, and regeneration as an alternative to
static routing.

The simulation results confirm the importance of temporal awareness and
resource recovery mechanisms for maintaining stable quantum communication
paths. The findings demonstrate that future quantum networks require routing
strategies capable of adapting to continuously changing entanglement states.

Overall, this work contributes toward the development of reliable and
resource-efficient routing techniques for future large-scale quantum
network infrastructures.